% Chapitre 2 - Page 1 | Travaux Avant-Projet | L3 ELN
% Version mise a jour avec les logos fournis par l'enseignant.
\documentclass[11pt,a4paper]{article}
\usepackage{fontspec}
\setmainfont{TeXGyrePagella}
\usepackage[french]{babel}
\usepackage{tikz}
\usepackage{graphicx}
\usetikzlibrary{arrows.meta,calc,positioning}
\usepackage{xcolor}
\usepackage{geometry}
\geometry{margin=0pt}
\pagestyle{empty}
\definecolor{bleufonce}{HTML}{1B3A6B}
\definecolor{bleupale}{HTML}{D9E9FB}
\definecolor{bleuclaire}{HTML}{EAF3FB}
\definecolor{protblue}{HTML}{2D63A7}
\definecolor{tinkblue}{HTML}{1479C9}
\definecolor{tinkyellow}{HTML}{F7C531}
\definecolor{tinkorange}{HTML}{F38B2C}
\definecolor{tinkred}{HTML}{DF4D44}
\definecolor{tinkgreen}{HTML}{7EC242}
\definecolor{fritgreen}{HTML}{7BC043}
\definecolor{fritorange}{HTML}{F39C12}
\definecolor{fritturq}{HTML}{26A69A}
\newcommand{\pos}[4]{\node[anchor=north west,inner sep=0pt,outer sep=0pt,text width=#3,align=left] at ($(current page.north west)+(#1,-#2)$) {#4};}
\newcommand{\sect}[1]{{\color{bleufonce}\bfseries\fontsize{17}{19}\selectfont #1}}
\newcommand{\bodyfont}{\fontsize{10.5}{13.2}\selectfont}
\begin{document}
\begin{tikzpicture}[remember picture,overlay]
  \path[use as bounding box] (0,0) rectangle (0,0);
  % Bordure extérieure
  \draw[bleufonce,line width=.95pt,rounded corners=8mm]
     ($(current page.north west)+(5mm,-5mm)$) rectangle
     ($(current page.south east)+(-5mm,5mm)$);
  % En-tête
  \node[anchor=north,align=center,inner sep=0pt] at ($(current page.north)+(0,-15.2mm)$)
    {{\color{bleufonce}\bfseries\fontsize{19}{21}\selectfont Cours : Travaux Avant-Projet}\\[-.5mm]
     {\color{bleufonce!85}\fontsize{10.5}{12}\selectfont rédigé par M REGANI Said}};
  \node[anchor=north,align=center,text width=181mm,inner sep=0pt]
   at ($(current page.north)+(0,-36.8mm)$)
    {{\color{bleufonce}\bfseries\fontsize{17.5}{21}\selectfont Chapitre 2 --- Logiciels de conception\\ et de simulation électronique}};

  % Introduction
  \pos{10mm}{62.5mm}{188mm}{\sect{Introduction}}
  \pos{10mm}{74mm}{187mm}{{\bodyfont Les logiciels de conception électronique permettent de dessiner un circuit, d'étudier son fonctionnement et de préparer sa réalisation. Ils aident à détecter certaines erreurs avant de passer au montage réel.}}

  % Objectifs pédagogiques
  \pos{10mm}{96mm}{188mm}{\sect{Objectifs pédagogiques}}
  \pos{10mm}{108.5mm}{187mm}{{\bodyfont À la fin de ce chapitre, l'étudiant sera capable de :}}
  \pos{16mm}{117mm}{181mm}{{\bodyfont
      \begin{minipage}{\linewidth}
      \setlength{\parindent}{0pt}\setlength{\parskip}{2.1mm}
      \textcolor{bleufonce}{\textbullet}\enspace Identifier le rôle des logiciels Proteus, Tinkercad et Fritzing.\\[1.7mm]
      \textcolor{bleufonce}{\textbullet}\enspace Réaliser et simuler un circuit électronique simple avec Proteus.\\[1.7mm]
      \textcolor{bleufonce}{\textbullet}\enspace Monter et simuler un circuit Arduino avec Tinkercad Circuits.\\[1.7mm]
      \textcolor{bleufonce}{\textbullet}\enspace Représenter un câblage et découvrir le dessin d'un PCB avec Fritzing.
      \end{minipage}}}

  % Encadré exemple
  \draw[bleufonce,line width=.8pt,rounded corners=2.8mm]
     ($(current page.north west)+(10mm,-151mm)$) rectangle
     ($(current page.north west)+(200mm,-182mm)$);
  \pos{16mm}{156mm}{174mm}{{\color{bleufonce}\bfseries\fontsize{12}{14}\selectfont Exemple fil rouge}}
  \pos{16mm}{166mm}{175mm}{{\bodyfont La \emph{lampe automatique} du chapitre 1 servira d'exemple pour étudier trois approches : simulation d'un circuit, montage Arduino virtuel et représentation du câblage.}}

  % Trois logiciels
  \pos{10mm}{188mm}{188mm}{\sect{Les trois logiciels du chapitre}}
  \pos{10mm}{200mm}{187mm}{{\bodyfont Les trois logiciels ont des fonctions spécifiques ; ils ne suivent aucun ordre imposé :}}
  \pos{16mm}{209mm}{181mm}{{\bodyfont
    \textcolor{bleufonce}{\textbullet}\enspace \textbf{Proteus :} dessiner un schéma électrique, simuler le circuit et observer les résultats.\\[2.5mm]
    \textcolor{bleufonce}{\textbullet}\enspace \textbf{Tinkercad Circuits :} réaliser un montage virtuel et simuler la programmation d'Arduino.\\[2.5mm]
    \textcolor{bleufonce}{\textbullet}\enspace \textbf{Fritzing :} représenter le câblage, réaliser un schéma et préparer un PCB ; ce n'est pas un simulateur de circuit.
  }}

  % Logos originaux fournis par l'enseignant. Chaque visuel reste dans sa carte
  % avec une zone distincte pour la legende, afin d'eviter tout chevauchement.
  % Carte Proteus : pictogramme original
  \begin{scope}[shift={($(current page.north west)+(15mm,-271mm)$)},x=1mm,y=1mm]
    \draw[bleufonce,line width=.7pt,rounded corners=2mm,fill=bleupale!40]
      (0,0) rectangle (54,25);
    \node[inner sep=0pt] at (27,12.5)
      {\includegraphics[width=23mm,height=17mm,keepaspectratio]{chapitre_2_logo_assets/logo_proteus.png}};
  \end{scope}

  % Carte Tinkercad : logotype fourni (fond transparent)
  \begin{scope}[shift={($(current page.north west)+(78mm,-271mm)$)},x=1mm,y=1mm]
    \draw[bleufonce,line width=.7pt,rounded corners=2mm,fill=bleupale!40]
      (0,0) rectangle (54,25);
    \node[inner sep=0pt] at (27,12.5)
      {\includegraphics[width=46mm,height=14mm,keepaspectratio]{chapitre_2_logo_assets/logo_tinkercad.png}};
  \end{scope}

  % Carte Fritzing : icone fournie
  \begin{scope}[shift={($(current page.north west)+(141mm,-271mm)$)},x=1mm,y=1mm]
    \draw[bleufonce,line width=.7pt,rounded corners=2mm,fill=bleupale!40]
      (0,0) rectangle (54,25);
    \node[inner sep=0pt] at (27,12.5)
      {\includegraphics[width=20mm,height=20mm,keepaspectratio]{chapitre_2_logo_assets/logo_fritzing.png}};
  \end{scope}

  \pos{10mm}{279mm}{189mm}{{\fontsize{9.4}{11.2}\selectfont Le choix du logiciel dépend du besoin du projet ; la simulation ne remplace pas les essais réels.}}

  % Pied de page
  \draw[bleufonce!55,line width=.45pt] ($(current page.south west)+(10mm,14.7mm)$) --
                                     ($(current page.south east)+(-10mm,14.7mm)$);
  \node[anchor=south west,inner sep=0pt,text=bleufonce!80] at ($(current page.south west)+(10mm,7.4mm)$)
    {\fontsize{9.5}{11}\selectfont 2026/2027};
  \node[anchor=south east,inner sep=0pt,text=bleufonce!80] at ($(current page.south east)+(-10mm,7.4mm)$)
    {\fontsize{9.5}{11}\selectfont Page 1};
\end{tikzpicture}
\null
\end{document}
